\documentclass[../../main.tex]{subfiles}
\begin{document}

\begin{center}
  {\fontsize{48pt}{58pt}\selectfont \textbf{1979}\normalsize\textbf{年度}}
\end{center}

\begin{table}[htb]
  \centering
  \caption{出題テーマと難易度}
  \begin{tabular}{|c|p{9cm}|c|}
    \hline
    \textbf{問題番号} & \textbf{テーマ} & \textbf{難易度} \\
    \hline
    第1問 & 定積分で直交するモニック多項式 & \\
    \hline
    第2問 & 三角関数を含む恒等式の係数決定 & \\
    \hline
    第3問 & 三角形の成立条件と関数の極値 & \\
    \hline
    第4問 & 偏りのあるサイコロの引き分け確率 & \\
    \hline
    第5問 & 円に内接する三角形の垂心と重心 & \\
    \hline
    第6問 & 指数的に縮む回転運動の軌跡 & \\
    \hline
  \end{tabular}
\end{table}

\begin{table}[htb]
  \centering
  \caption{解答}
  \renewcommand{\arraystretch}{1.4}
  \begin{tabular}{|c|c|p{9cm}|}
    \hline
    \textbf{問題番号} & \textbf{小問} & \textbf{答え} \\
    \hline
    第1問 & (1) & $P_1(x)=x$ \\
    \cline{2-3}
    & (2) & $P_2(x)=x^2-\dfrac{1}{3}$ \\
    \cline{2-3}
    & (3) & $P_3(x)=x^3-\dfrac{3}{5}x$ \\
    \hline
    第2問 & -- & $a=b=\dfrac{1}{2}$，$c=\pi+2n\pi$（$n\in\mathbb{Z}$） \\
    \hline
    第3問 & (1) & $\min f=2+\sqrt{3}$，$\max g=2-\sqrt{3}$ \\
    \cline{2-3}
    & (2) & $2-\sqrt{3}<p<2+\sqrt{3}$ \\
    \hline
    第4問 & (1) & $P=\sum_{k=1}^{6}p_k^2$ \\
    \cline{2-3}
    & (2) & $P\ge\dfrac{1}{6}$．等号はすべての$p_k=\dfrac{1}{6}$のとき \\
    \hline
    第5問 & (1) & $\overrightarrow{OH}=\sum_{i=1}^{3}\overrightarrow{OA_i}$ \\
    \cline{2-3}
    & (2) & 共通の定点$X$は$\overrightarrow{OX}=\dfrac{1}{4}\sum_{i=1}^{6}\overrightarrow{OA_i}$ \\
    \hline
    第6問 & -- &  \\
    \hline
  \end{tabular}
\end{table}

\paragraph{試験形態}
6問．試験時間は要確認．

\end{document}
